\clearpage
\section{The \textsf{reverse-copy} and \textsf{reverse} Algorithms}
\label{reversecopy}

\subsection{\textsf{reverse-copy}}

We now study the \texttt{reverse-copy} algorithm from ``ACSL By Example''.
This algorithm copies the contents from a source sequence to a destination
sequence in reverse order.

The specification of the \texttt{reverse-copy} algorithm is as follows:

\listingname{reversecopy.h}
\cinput{reversecopy/reversecopy.spec}

The implementation of the algorithm is:

\listingname{reversecopy.c}
\cinput{reversecopy/reversecopy.impl}

Running the \textsf{WP} plug-in does not allow proving all properties:
\fclogs{mutating}{reversecopy}

\subsection{\textsf{reverse}}

The \texttt{reverse} algorithm is an in place version of the
\texttt{reverse-copy} algorithm.

Its specification is as follows:

\listingname{reverse.h}
\cinput{reversecopy/reverse.spec}

The implementation is:

\listingname{reverse.c}
\cinput{reversecopy/reverse.impl}

Running the \textsf{WP} plug-in does not allow proving some loop invariants
preservation nor certain loop assigns:
\fclogs{mutating}{reverse}
